% 本文件的姓名、联系方式、学校、公司与经历均为虚构示例。
% 请在自己的私有副本中替换内容；发布前检查 PDF、图片和元数据。
% 特殊字符：百分号写作 \%，下划线写作 \_，& 写作 \&。
\ResumeHeader{张三}{2002.06}{浙江杭州}{群众}{男}{13800000000}{hello@example.com}

\ResumeSection{0}{教育背景}
\vspace{-3bp}
\Education{2024.09\hspace{.3em}–\hspace{.3em}2027.06}{示例大学 计算机学院}{计算机技术（硕士）}
\Education{2020.09\hspace{.3em}–\hspace{.3em}2024.06}{示例大学 信息学院}{软件工程（本科）}
\begin{resumeitems}
  \Point{主修课程：}{数据结构、操作系统、数据库系统、机器学习。}
  \Point{学习经历：}{参与课程项目与实验室讨论，关注软件工程和人工智能应用。}
  \item 请根据真实情况填写成绩、荣誉和学习经历，避免使用未经核实的数据。
\end{resumeitems}

\ResumeSection{1}{实习经历}
\Entry{某科技公司｜软件开发实习生}{2023.07 -{}- 2023.10}
\begin{resumeitems}
  \Point{需求与实现：}{参与内部工具的需求梳理，完成数据查询页面和接口开发。}
  \Point{测试与协作：}{编写单元测试和使用说明，与团队共同排查问题、改进交付流程。}
\end{resumeitems}

\ResumeSection{2}{项目经历}
\Entry{校园活动管理平台}{2024.03 -{}- 2024.06}
\Description{面向学生社团的课程示例项目，支持活动发布、报名管理和结果导出。}
\begin{resumeitems}
  \Point{系统设计：}{梳理角色权限与数据模型，设计活动、报名和通知模块。}
  \Point{功能开发：}{实现列表筛选、表单校验和批量导出，并完善错误提示。}
  \Point{项目复盘：}{编写部署文档，记录测试方法与后续可改进的功能。}
\end{resumeitems}
\Entry{文本分类实验}{2023.09 -{}- 2023.12}
\Description{在公开数据集上比较基线模型与深度学习模型的分类表现。}
\begin{resumeitems}
  \Point{数据处理：}{完成清洗、划分和类别分布检查，避免训练与验证数据混用。}
  \Point{模型评估：}{使用准确率、召回率和混淆矩阵分析模型的适用场景。}
  \Point{成果展示：}{整理实验过程、复现步骤和可视化结果。}
\end{resumeitems}

\ResumeSection{3}{荣誉奖项}
\AwardRow{校级课程设计展示奖（示例）}{院级志愿服务表彰（示例）}
\AwardRow{项目实践优秀报告（示例）}{请替换为自己的真实奖项}

\vspace{4bp}
\ResumeSection{4}{技能特长}
\SkillHead{开发与数据能力}{协作与工程能力}
\SkillRow{熟悉 Python、SQL 和常用数据处理工具}{能够使用 Git 管理代码与协作}
\SkillRow{了解常见机器学习模型与评估方法}{熟悉 Linux 基本命令与部署流程}
\SkillRow{能够编写清晰的技术文档和测试用例}{具备问题定位与跨团队沟通能力}
